\scanpage{82}

§. \uline{Inégalités de convexité et degrés remarquables d'opérateurs compacts}

\deleted{\unsure{La méthode de convexité}} \inserted{\textcircled{\scriptsize Th. 1}} \textcircled{1°} \deleted{\uline{Lemme fondamental}} \inserted{(H. Weyl).} Considérons deux suites finies

$(\alpha_i)_{1\leq i\leq n}$ et $(\beta_i)_{\unsure{1\leq i\leq n|i\leq n}}$ de nbres réels \inserted{on suppose} \deleted{la deuxième \illegible{1} \unsure{décroissante}}

\[
\boxed{\alpha_1\geq\alpha_2\ldots\geq\alpha_n}\quad\text{et}\quad
\boxed{\alpha_1+\cdots+\alpha_{\unsure{M|N|m}}\leq\beta_1+\cdots+\beta_M\enspace\text{pour }M=1,\ldots n.}\tag*{(1)\quad(2)}
\]

Soit $\varphi(t_1,\ldots,t_n)$ une fonction numérique \uline{convexe} et \uline{symétrique} sur $R^n$. On suppose vérifiée l'une des deux conditions suivantes : a) $\varphi$ est croissante \deleted{\unsure{(pour}}\inserted{en} chaque argument $t_i$ \deleted{la restriction d'un ensemble sur $R^n$} b) $\sum_{i=1}^n\alpha_i=\sum_{i=1}^{n}\beta_i$.

% slanted
\margin{\unsure{fonction} où l'on \unsure{devrait} \unsure{admettre} \unsure{H} \illegible{1} \deleted{\unsure{partie}} \unsure{un} convexe ouvert.}

Sous ces conditions, on a

\[
\varphi(\alpha_1,\ldots,\alpha_n)\leq\varphi(\beta_1,\ldots\beta_n)\tag*{(3)}
\]

\deleted{Appliquons \unsure{tout de suite} le corollaire suivant, qui \unsure{seul nous} \illegible{1} \unsure{pour la} \unsure{question} qui \unsure{effacerait} \illegible{1}. On en conclut immédiatement le}

\deleted{\uline{Corollaire 1}} \inserted{\textcircled{2°}} Soit $f(\rho_1,\ldots,\rho_n)$ fonction \deleted{\unsure{continue|numérique}} \inserted{\uline{numérique symétrique}} de $n$ arguments \deleted{\unsure{positifs}} \inserted{$\rho_i>0$}, telle que $f(e^{t_1},\ldots,e^{t_n})$ soit une fonction \uline{convexe} de $(t_1,\ldots t_n)\in R^n$. Soient $(\rho_1,\ldots\rho_n)$ et $(\sigma_1,\ldots,\sigma_n)$ deux suites \inserted{finies} positives, telles que \deleted{\unsure{(Hies.)}} \textcircled{4} $\boxed{\rho_1\geq\rho_2\ldots\geq\rho_n}$ et que \deleted{\unsure{\ensuremath{\sum\rho_i}}} \textcircled{5} $\boxed{\rho_1\cdots\rho_M\leq\sigma_1\cdots\sigma_M}$ pour \deleted{tout} $M=1,\ldots n$. On suppose vérifiée l'une des deux conditions suivantes :

% slanted
\margin{Il suffit que l'ensemble où est définie $f$ soit une partie \unsure{logarithmiquement} \unsure{convexe} \unsure{convenable} et symétrique.}

a) $f$ est \inserted{définie et continue pour $\rho_i\geq0$, et} fonction croissante \deleted{des} \deleted{\unsure{chacun des|chaque argument $\rho_i$}}\inserted{de chaque argument} $\rho_i$ b) $\rho_1\cdots\rho_n=\sigma_1\cdots\sigma_n\ne0$. Sous ces conditions, on a

\[
f(\rho_1,\ldots\rho_n)\leq f(\sigma_1,\ldots\sigma_n)\tag*{\deleted{\unsure{(4)}} (6)}
\]

En effet, si les $\rho_i,\sigma_i$ sont $>0$, cet énoncé est manifestement équivalent à l'énoncé du \deleted{lemme}\inserted{th.}, en posant $\alpha_i=\log\rho_i$, $\beta_i=\log\sigma_i$. \inserted{Cela prouve le \deleted{\unsure{corollaire}}\inserted{(2°)} en particulier dans le cas b).}

\deleted{\unsure{On en conclut}, puisque} \inserted{Montrons que} dans le cas a), tout \deleted{l'un et des} systèmes $(\rho_i),(\sigma_i)$ tels que $\rho_i\geq0$, $\deleted{\unsure{\sigma_i\geq0|\rho_i>0}}$ \inserted{$\rho_1\geq\rho_2\ldots\geq\rho_n\geq0$,} $\rho_1\cdots\rho_M\leq\sigma_1\cdots\sigma_M$ pour $M=1,\ldots n$, \deleted{\unsure{et naturellement} $\rho_1\ldots\rho_n=\sigma_1\ldots\sigma_n$)} \deleted{\unsure{et l'adhérence}} \inserted{est limite dans $R^{2n}$} \deleted{\unsure{de sa partie formée des systèmes tels}} \inserted{de systèmes analogues} $(\rho'_i,\sigma'_i)$ avec $\rho'_i>0$, $\sigma'_i>0$ (d'où résultera le \deleted{corollaire}\inserted{(2°) dans le cas a)} par prolongement par continuité). En effet, si \deleted{\unsure{les $\rho_i$}} \inserted{$\sigma_i>0$ pour tout $i$,} on voit qu'en \deleted{\unsure{tous les cas}} augmentant \inserted{de suffisamment peu} les $\rho_i$ qui seraient nuls, \deleted{\unsure{sont nuls, et il suffit de reprendre les inégalités pour|et il suffit d'approcher les systèmes}} on satisfait encore aux conditions imposées, \deleted{\unsure{et si} les systèmes $(\rho'_i)$ et $(\sigma'_i)$, avec $\rho'_i=\sigma'_i=\frac1n$ pour tout $i$. D'autre part} d'où la conclusion dans ce cas.

\margin{\unsure{?}}

\deleted{Si plus généralement, \unsure{dans|donc} \illegible{1} \unsure{l'on} \illegible{2} \unsure{qui sont} \illegible{1} d'abord \illegible{1}}

% archive sheet number: 83
